\selectlanguage{american}%

\section{Conjugate Gradient}

Consider any algorithm that starts with $x^{(0)}=0$ and satisfies
the invariant
\[
x^{(k+1)}\in x^{(0)}+\spn\{\nabla f(x^{(0)}),\nabla f(x^{(1)}),\cdots,\nabla f(x^{(k)})\}.
\]
For $f(x)=\frac{1}{2}x^{\top}Ax-b^{\top}x$, it is easy to see that
$x^{(k)}\in\mathcal{K}_{k}$ defined as follows:
\begin{defn}
Define the Krylov subspaces $\K_{0}=\{0\}$, $\K_{k}=\spn\{b,Ab,\cdots,A^{k-1}b\}$.
\end{defn}

Therefore, the best such algorithm can do is to solve $\min_{x\in\K_{k}}f(x)$.
For the quadratic function $f(x)\defeq\frac{1}{2}x^{\top}Ax-b^{\top}x$,
one can compute $\min_{x\in\K_{k}}f(x)$ efficiently using the conjugate
gradient method. Note that 
\[
\arg\min_{x\in\K_{k}}f(x)=\arg\min_{x\in\K_{k}}\norm{x-x^{*}}^{2}_{A}
\]
where $Ax^{*}=b$.

Throughout this section, we assume $A$ is positive definite. For
a consistent full-column-rank linear system $Ax=b$ (or for the associated
least-squares problem), we can pass to $A^{\top}Ax=A^{\top}b$ to
satisfy the assumption. This transformation may square the condition
number, so it is mainly a reduction to the positive definite setting
rather than a numerically preferable method.
\begin{lem}
Let $x^{(k)}=\argmin_{x\in\K_{k}}f(x)$ be the Krylov sequence. Then,
the steps $v^{(k)}=x^{(k)}-x^{(k-1)}$ are ``conjugate'', namely,
\[
v^{(i)\top}Av^{(j)}=0\text{ for all }i\neq j.
\]
\end{lem}

\begin{proof}
Assume that $i<j$. The optimality of $x^{(j)}$ shows that $\nabla f(x^{(j)})\in\K^{\perp}_{j}\subset\K^{\perp}_{j-1}$
and that $\nabla f(x^{(j-1)})\in\K^{\perp}_{j-1}$. Hence, we have
\[
Av^{(j)}=\nabla f(x^{(j)})-\nabla f(x^{(j-1)})\in\K^{\perp}_{j-1}.
\]
Next, we note that $v^{(i)}=x^{(i)}-x^{(i-1)}\in\mathcal{K}_{i}\subset\mathcal{K}_{j-1}$.
Hence, we have $v^{(i)\top}Av^{(j)}=0$.
\end{proof}
Since the steps are conjugate, the vectors $v^{(i)}$ form a conjugate
basis for the Krylov subspaces:
\[
\K_{k}=\spn\{v^{(1)},v^{(2)},\cdots,v^{(k)}\}.
\]

Note that $x^{(k)}=x^{(k-1)}+v^{(k)}$. Hence, it suffices to find
a formula for $v^{(k)}$.
\begin{lem}
\label{lem:CG_v}For $k\ge2$, we have 
\[
v^{(k)}=\frac{v^{(k)\top}r^{(k-1)}}{\|r^{(k-1)}\|^{2}}\left(r^{(k-1)}-\frac{r^{(k-1)\top}Av^{(k-1)}}{v^{(k-1)\top}Av^{(k-1)}}v^{(k-1)}\right).
\]
\end{lem}

\begin{proof}
Again, the optimality condition shows that $\nabla f(x^{(k-1)})\in\K^{\perp}_{k-1}$
and that $\nabla f(x^{(k-1)})\in\K_{k}$ by the definition of $\K$.
Hence, we have
\[
\K_{k}=\spn\{v^{(1)},\cdots,v^{(k-1)},r^{(k-1)}\}
\]
where $r^{(k-1)}=b-Ax^{(k-1)}$. Therefore, we have
\[
v^{(k)}=c_{0}r^{(k-1)}+\sum^{k-1}_{i=1}c_{i}v^{(i)}
\]
for some $c_{0}$. 

For $c_{0}$, we use that $r^{(k-1)}\in\K^{\perp}_{k-1}$. This gives
$v^{(k)\top}r^{(k-1)}=c_{0}\|r^{(k-1)}\|^{2}$.

For $c_{k-1}$, since $v^{(i)\top}Av^{(k-1)}=0$ for any $i\neq k-1$,
we have
\[
0=v^{(k)\top}Av^{(k-1)}=c_{0}r^{(k-1)\top}Av^{(k-1)}+c_{k-1}v^{(k-1)\top}Av^{(k-1)}
\]

and $c_{k-1}=-c_{0}\cdot\frac{r^{(k-1)\top}Av^{(k-1)}}{v^{(k-1)\top}Av^{(k-1)}}$. 

For other $c_{i}$, we note that $Av^{(j)}\in\K_{j+1}$ and $r^{(k-1)}\in\K^{\perp}_{k-1}$.
Hence, $c_{i}=0$ for all $i\notin\{0,k-1\}$. 

Substituting the $c_{i}$, we get
\[
v^{(k)}=\frac{v^{(k)\top}r^{(k-1)}}{\|r^{(k-1)}\|^{2}}\left(r^{(k-1)}-\frac{r^{(k-1)\top}Av^{(k-1)}}{v^{(k-1)\top}Av^{(k-1)}}v^{(k-1)}\right).
\]
\end{proof}
To make the formula simpler, we define $p^{(k-1)}=\frac{\|r^{(k-1)}\|^{2}_{2}}{v^{(k)\top}r^{(k-1)}}v^{(k)}$
for $k\ge1$.
\begin{lem}
We have that $x^{(k)}=x^{(k-1)}+\frac{\|r^{(k-1)}\|^{2}_{2}}{p^{(k-1)\top}Ap^{(k-1)}}p^{(k-1)}$
and, for $k\ge2$, $p^{(k-1)}=r^{(k-1)}+\frac{\|r^{(k-1)}\|^{2}}{\|r^{(k-2)}\|^{2}_{2}}p^{(k-2)}.$
\end{lem}

\begin{proof}
For the first formula, note that
\[
x^{(k)}=x^{(k-1)}+v^{(k)}=x^{(k-1)}+\frac{v^{(k)\top}r^{(k-1)}}{\|r^{(k-1)}\|^{2}_{2}}p^{(k-1)}.
\]
For the quantity $v^{(k)\top}r^{(k-1)}$, we note that $f(x^{(k-1)}+tv^{(k)})$
is minimized at $t=1$ and hence
\begin{align*}
v^{(k)\top}r^{(k-1)} & =v^{(k)\top}Av^{(k)}\\
 & =\frac{v^{(k)\top}r^{(k-1)}}{\|r^{(k-1)}\|^{2}_{2}}\frac{v^{(k)\top}r^{(k-1)}}{\|r^{(k-1)}\|^{2}_{2}}p^{(k-1)\top}Ap^{(k-1)}
\end{align*}
where we used the definition of $p^{(k-1)}$. Simplifying gives
\[
\frac{v^{(k)\top}r^{(k-1)}}{\|r^{(k-1)}\|^{2}_{2}}=\frac{\|r^{(k-1)}\|^{2}_{2}}{p^{(k-1)\top}Ap^{(k-1)}}
\]
and hence the first formula.

For the second formula, we use Lemma \ref{lem:CG_v} and substitute
$p^{(k-1)}=\frac{\|r^{(k-1)}\|^{2}_{2}}{v^{(k)\top}r^{(k-1)}}v^{(k)}$.
This gives
\begin{equation}
p^{(k-1)}=r^{(k-1)}-\frac{r^{(k-1)\top}Av^{(k-1)}}{\|r^{(k-2)}\|^{2}_{2}}p^{(k-2)}.\label{eq:CG_p_temp}
\end{equation}
Note that
\[
r^{(k-1)}=b-Ax^{(k-1)}=r^{(k-2)}-Av^{(k-1)}.
\]
Taking inner product with $r^{(k-1)}$ and using $r^{(k-1)}\perp r^{(k-2)}$
gives $\|r^{(k-1)}\|^{2}=-r^{(k-1)\top}Av^{(k-1)}$. Putting this
into (\ref{eq:CG_p_temp}) gives the result.
\end{proof}
We use the standard zero-indexed search directions: $r^{(0)}=b-Ax^{(0)}$,
$p^{(0)}=r^{(0)}$, and iteration $k$ produces $x^{(k+1)}$, $r^{(k+1)}$,
and $p^{(k+1)}$.

\begin{algorithm2e}[H]

\caption{$\mathtt{ConjugateGradient}$ (CG)}

\SetAlgoLined

$r^{(0)}=b-Ax^{(0)}$. $p^{(0)}=r^{(0)}$.

\For{$k=0,1,2,\cdots$}{

$\alpha_{k}=\frac{\|r^{(k)}\|^{2}_{2}}{p^{(k)\top}Ap^{(k)}}$.

$x^{(k+1)}\leftarrow x^{(k)}+\alpha_{k}p^{(k)}$.

$r^{(k+1)}\leftarrow r^{(k)}-\alpha_{k}Ap^{(k)}$.

\lIf{$\|r^{(k+1)}\|_{2}\leq\epsilon$}{\Return$x^{(k+1)}$}

$\beta_{k}=\frac{\|r^{(k+1)}\|^{2}}{\|r^{(k)}\|^{2}}$.

$p^{(k+1)}=r^{(k+1)}+\beta_{k}p^{(k)}.$

}

\end{algorithm2e}
\begin{thm}
\label{thm:CG}Let $f(x)=\frac{1}{2}x^{\top}Ax-b^{\top}x$ for some
positive definite matrix $A$ and some vector $b$. Let $x^{(k)}$
be the sequence produced by $\mathtt{ConjugateGradient}$. Then, we
have
\[
f(x^{(k)})-f(x^{*})\leq\frac{1}{2}\inf_{\deg q\leq k,\,q(0)=1}\max_{i}q(\lambda_{i})^{2}\cdot\|b\|^{2}_{A^{-1}}
\]
where $q$ ranges over polynomials of degree at most $k$ and $\lambda_{i}$
are eigenvalues of $A$.
\end{thm}

\begin{proof}
Note that $\K_{k}=\{p(A)b:\deg p<k\}$. Hence, we have
\begin{align*}
2(f(x^{(k)})-f(x^{*})) & =\min_{x\in\K_{k}}\|x-x^{*}\|^{2}_{A}\\
 & =\inf_{\deg p<k}\|(p(A)-A^{-1})b\|^{2}_{A}\\
 & =\inf_{\deg p<k}\|A^{-\frac{1}{2}}(Ap(A)-I)A^{-\frac{1}{2}}b\|^{2}_{A}\\
 & =\inf_{\deg p<k}\|(Ap(A)-I)A^{-\frac{1}{2}}b\|^{2}_{2}
\end{align*}
Let $q(x)=1-xp(x)$. Note that $q(0)=1$ and $\deg q\leq k$ and any
such $q$ is of the form $1-xp$. Hence, we have
\begin{align*}
2(f(x^{(k)})-f(x^{*})) & =\inf_{\deg q\leq k,q(0)=1}\|q(A)A^{-\frac{1}{2}}b\|^{2}_{2}\\
 & \leq\inf_{\deg q\leq k,q(0)=1}\|q(A)\|^{2}_{\op}\cdot\|b\|^{2}_{A^{-1}}.
\end{align*}
The result follows from the fact that $\|q(A)\|^{2}_{\op}=\max_{i}q(\lambda_{i})^{2}$.
\end{proof}
Assume that $0<\mu$ and $\mu I\preceq A\preceq LI$. It is known
that there is a polynomial $q$ of degree $k$ with $q(0)=1$ such
that
\[
\max_{i}q(\lambda_{i})^{2}\leq2\left(1-\frac{2}{\sqrt{\frac{L}{\mu}}+1}\right)^{k}
\]
Therefore, it takes $O(\sqrt{\frac{L}{\mu}}\log(\frac{1}{\epsilon}))$
iterations to find $x$ such that $f(x)-f(x^{*})\leq\epsilon\cdot\|b\|^{2}_{A^{-1}}$.

Also, we note that if there are only $s$ distinct eigenvalues in
$A$, then conjugate gradient finds the exact solution in $s$ iterations.

\selectlanguage{english}%

